\clearpage
\subsection{置信度阈值与点云融合结果}
\subsubsection{点云形状与保留范围}

使用同一组 5 视图预测，分别设置 0.6、0.8 和 0.95 三个置信度阈值，保持几何检查和体素大小不变。三组点云采用相同观察方向、范围和点大小。

\ReportFigure[173bp]{threshold_clouds.png}{5 视图组的点云融合结果。置信度阈值逐渐提高后，建筑周围的零散点减少，部分表面也变得稀疏。}

\ReportFigure[143bp]{threshold_masks.png}{参考视角 0012 的保留区域。浅灰色为被过滤的像素，其余位置保留原图颜色；三图采用同一参考图和像素范围。}

\subsubsection{点数与重投影偏差比较}
\input{threshold-table.tex}

\ReportNote{“保留观测点数”为过滤后、体素合并前的总数；同一表面可能被多个视角重复观测。“重投影偏差”先对每个参考像素计算其在有效源图中的往返像素偏差均值，再对保留像素取中位数，最后对 49 个视角求平均。统计包含能完成投影的源图，不只包含通过阈值的源图；它仍属于内部一致性指标。}

阈值从 0.6 提高到 0.95 后，保留比例由 66.96\% 降至 55.91\%，体素合并后的点数由约 106.7 万降至 84.4 万，减少约 21.0\%。重投影偏差也由 0.245 px 降至 0.209 px，说明被保留的观测在本组图像之间更一致。\textbf{这种下降与筛选范围变小有关，不能直接解释成真实深度误差同样降低。}

从图中可以看到，三组都保留了建筑主体，提高阈值主要影响建筑周围的零散点和部分表面细节。阈值 0.95 的结果较干净，但参考图中的保留区域也缩小了。0.8 的结果在过滤杂点的同时保留了较多观测，因此我们将它作为本次实验的默认设置；对于完整性更重要的场景，则需要结合缺失区域重新选择阈值。

